This section separates what the descriptive pass suggests from what the inferential stage decides. The first is reported here only to establish that there is something to measure; the second is reported under the rules fixed in \autoref{tab:bayesDecisions}, and returns no verdict it has not earned.

\begin{table}[!t]
    \centering
    \small
    \begin{tabular}{@{}cccc@{}}
        \toprule
        \textbf{$P$} & \textbf{$S$} &
        $|\mathcal{F}_{S}\setminus\mathcal{F}_{P}|$ & $|\mathcal{F}_{P}\setminus\mathcal{F}_{S}|$ \\
        \midrule
        $8$ & $8$ & $0$ & $0$ \\
        \midrule
        $16$ & $8$ & $0$ & $4$ \\
        $16$ & $12$ & $4$ & $6$ \\
        $16$ & $16$ & $0$ & $0$ \\
        \midrule
        $24$ & $8$ & $0$ & $8$ \\
        $24$ & $12$ & $0$ & $6$ \\
        $24$ & $16$ & $4$ & $8$ \\
        $24$ & $20$ & $8$ & $10$ \\
        $24$ & $24$ & $0$ & $0$ \\
        \midrule
        $32$ & $8$ & $0$ & $12$ \\
        $32$ & $12$ & $4$ & $14$ \\
        $32$ & $16$ & $0$ & $8$ \\
        $32$ & $20$ & $8$ & $14$ \\
        $32$ & $24$ & $8$ & $12$ \\
        $32$ & $32$ & $0$ & $0$ \\
        \midrule
        \multicolumn{2}{c}{\textbf{Total}} & $36$ & $102$ \\
        \bottomrule
    \end{tabular}
\caption{What each geometry can attribute to one branch alone: the frequencies in $[2,250]$\,Hz that one grid predicts and the other does not. A zero in the first column is the case $S \mid P$, where every stride site $c f_s/S$ is also a patch site $k f_s/P$; a zero in both is $P=S$, where the two branches are the same set of frequencies. The last row totals them over the design. The counts are taken from the grids themselves rather than from any measurement, and are unchanged whether coincidence is required exactly or only to within the $1$\,Hz resolution of the sweep.}
    \label{tab:branchSites}
\end{table}


One limit is arithmetic rather than empirical, and it is worth stating before any estimate. The two branches of $\mathcal F_{\mathrm{lock}}$ are not equally well served by the design, because a geometry can only speak about the stride branch where the stride grid is not already contained in the patch grid. \autoref{tab:branchSites} counts, for each of the fifteen runs, the frequencies in the analysed band that one grid predicts and the other does not. Where $S \mid P$ every stride site is also a patch site, so nine of the fifteen runs yield no stride-only frequency at all, and four of those nine are the $P=S$ baselines where the two grids coincide entirely. The stride branch is therefore carried by the six runs with $S \nmid P$, and $\theta_S$ and $\kappa_S$ rest on those. Across the whole design they supply thirty-six stride-only frequencies against one hundred and two patch-only ones, and the most any single geometry supplies is eight, below the ten unambiguous sites that \autoref{sec:five} requires of a branch: the stride branch is identified only by pooling the six, never by any one of them. Those same six are the geometries whose analysis context is truncated, the ones with $L_{\mathrm{tok}}<480$ in \autoref{tab:hfModels}, so the subset that carries the stride branch is also the subset least like the rest of the design. All of this is fixed by the choice of $(P,S)$ before anything is measured, and it bounds what the stride-branch estimates below can be asked to settle.

Before the Bayesian estimates are reported, the pages of \autoref{sec:appSweep} provide a descriptive reading of what those models must explain. Across the retrained geometries, dominant-frequency reconstruction appears more reliable over the lower half of the analysed band, while faster inputs often return near $0$\,Hz. Within that lower half, possible holes occur near $\mathcal F_{\mathrm{lock}}$, and hierarchical band probes sometimes weaken there, although the local $f\pm1$\,Hz task remains stronger. The generated-signal examples add a second question: slow components may survive, yet their forecast amplitudes can approach the noise floor, making recovery ambiguous. These plots therefore motivate paired Bayesian analysis on generated backgrounds and frame the contrasts reported next. They are exploratory observations only, not confirmation of structural aliasing or of the preregistered hypotheses in \autoref{sec:three}.

\autoref{sec:appSolar} adds two descriptive readings of a different kind. It reads the probe by depth rather than pooled, so the stage at which a candidate frequency stops being linearly separable can be located inside the stack, which Model~B cannot report because it carries the stage as a nuisance offset by design; and it puts the models in front of the measured photovoltaic series, where the forecast error is of the order of the quantity predicted on the sunrise and afternoon ramps against about two per cent at midday. Both are descriptive: the representational claim is decided by $e^{\theta_{\mathrm{lock}}}$ under the rule of \autoref{tab:bayesDecisions} and nowhere else.
